\documentclass[11pt,a4paper]{article}
\usepackage[T1]{fontenc}\usepackage[margin=25mm]{geometry}
\usepackage{amsmath,amssymb,amsthm,mathtools,lmodern,microtype}
\usepackage[hidelinks]{hyperref}\usepackage{xurl}
\hypersetup{pdftitle={First Negative Degrees in Integer Thue Morse Pressure},pdfauthor={Research note prepared for Vladimir Reshetnikov with OpenAI}}
\setlength{\parskip}{2pt}
\newtheorem{theorem}{Theorem}[section]\newtheorem{lemma}[theorem]{Lemma}
\newtheorem{proposition}[theorem]{Proposition}
\newcommand{\ev}{\operatorname{ev}}\newcommand{\Z}{\mathbb Z}
\newcommand{\HH}{\mathcal H}\newcommand{\GG}{\mathcal G}
\title{First Negative Degrees in\\Integer Thue--Morse Pressure}
\author{Research note prepared for Vladimir Reshetnikov with OpenAI}
\date{1 October 2026}
\begin{document}\maketitle
\begin{abstract}
Let $N_m$ be the first even degree strictly above $2m$ at which the phase Taylor coefficient of integer Thue--Morse pressure is negative. We prove
$N_m=\gamma m-\Lambda^{-1}\log\log(2m)+O(1)$, where a positive-series saddle equation defines $\gamma\in(6.662966,6.662968)$ and $\Lambda>0$. In particular $N_m/m$ converges to a value strictly below $20/3$. The key comparison is between a positive first model and an actual quadratic Schur term with multiplier $2\log_2(2m)+O(1)$. A principal inverse-branch identity bounds the frozen correction at the second-model scale. Exact rational compact-domain certificates join the local saddle arcs to the rest of the coefficient contour. The result is asymptotic; it does not supply an unconditional floor rule or an effective onset threshold.
\end{abstract}
\paragraph{Status.} This is unrefereed ordinary mathematics with exact rational interval certificates, not a Lean formalization. Earlier analytic prerequisites are included unchanged with the sources. No exhaustive priority claim is made.

\section{Statement and the two positive models}
Use the cosine-phase normalization of the repository manuscript~\cite{Repo}, and put
\[
 P_m(t)=p_m(t/\pi),\qquad d=2m,\qquad
 N_m=\min\{N>d:N\text{ even},\ [t^N]P_m(t)<0\}.
\]
The infinite-sign theorem~\cite{Signs} ensures that $N_m$ exists for every integer $m\ge2$. The lower negative coefficients and the cancellation at degree $2m$ are excluded from this minimum.

Define
\[
 c=2/\pi,\quad b_j=\tan(\pi/2^{j+1}),\quad
 B(a)=\prod_{j\ge1}(1+ab_j),\quad L(a)=cB(a),\quad
 \GG(t)=\tan t\,L(\tan t).
\]
The product $B$ is entire. The series $\GG$ has strictly positive coefficients at every degree at least one and is analytic on $|t|<\pi/2$. Put
\[
 \mu(t)=\frac{t\GG'(t)}{\GG(t)}.
\]
For $\sigma>k$, let $t_k(\sigma)\in(0,\pi/2)$ solve $\mu(t_k)=\sigma/k$, and set
\begin{equation}\label{eq:rates}
 \Phi_k(\sigma)=k\log\GG(t_k)-\sigma\log t_k.
\end{equation}
The logarithmic derivative is the mean of a positive coefficient law; its derivative in $\log t$ is its strictly positive variance. Moreover $\mu$ increases from one to infinity, so each saddle exists and is unique.

There is a unique crossing $\sigma_*>2$ of $\Phi_1$ and $\Phi_2$. Indeed
\begin{equation}\label{eq:delta}
 \Delta(\sigma):=\Phi_2(\sigma)-\Phi_1(\sigma),\qquad
 \Delta'(\sigma)=\log(t_1(\sigma)/t_2(\sigma))>0,
\end{equation}
and the rational certificate gives opposite signs at $3.331483$ and $3.331484$. Write
\[
 \gamma=2\sigma_*,\qquad
 \Lambda=\log\frac{t_1(\sigma_*)}{t_2(\sigma_*)}.
\]
Numerical orientation gives $\gamma\approx6.6629666581$ and $\Lambda\approx0.8062663610$; the rigorous slope interval used here is the one above.
\begin{theorem}\label{thm:main}
As the integer $m$ tends to infinity,
\begin{equation}\label{eq:main}
 N_m=\gamma m-\frac{\log\log(2m)}{\Lambda}+O(1).
\end{equation}
In particular $N_m/m\to\gamma\in(6.662966,6.662968)$.
\end{theorem}
All unqualified logarithms are natural. The bounded term includes the even-degree lattice. No unconditional rounding formula is asserted.

\paragraph{Related context.}
Fan, Schmeling and Shen~\cite{FSS} study the same phase-dependent cosine-product family through ergodic optimization and multifractal analysis. Baake, Gohlke, Kesseb\"ohmer and Schindler~\cite{BGKS} study pressure and scaling for the classical Thue--Morse measure. The variable and asymptotic regime here concern high-order phase Taylor signs as the integer moment order grows. These references provide context, rather than an exhaustive novelty determination.

\section{Schur terms and the analytic prerequisites}
The normalized transfer operator is
\[
 (V_af)(x)=\sum_{2y=x\pmod1}(\cos\pi y+a\sin\pi y)^d f(y).
\]
Let $\lambda(a)$ be the eigenvalue branch of $V_a$ with $\lambda(0)=1$. The pressure germ is exactly
\[P_m(t)=d\log\cos t+\log\lambda(\tan t),\]
with the logarithm zero at the origin. Let $h_0$ be the fixed function of $V_0$, normalized by $h_0(0)=1$, let $\ev f=f(1/2)$, and put $Qf=f-f(0)h_0$. On functions vanishing at zero set
\begin{gather*}
 \mathcal B_a=QV_aQ,\qquad q_a=QV_ah_0,\qquad
 F=\ev h_0+\ev(I-\mathcal B_a)^{-1}q_a,\\
 G=\ev(I-\mathcal B_a)^{-2}q_a,\qquad
 H=\ev(I-\mathcal B_a)^{-3}q_a.
\end{gather*}
The functions are even analytic germs, with $G(0)=H(0)=0$. If $\delta=\lambda-1$ for the normalized simple eigenvalue, then
\[
 \delta=a^d\{F-\delta G+\delta^2H+O(\delta^3)\}.
\]
Consequently, below degree $4d$,
\begin{align}
 P_m(t)&=d\log\cos t+\tan^dt\,F(\tan t)
       -\tan^{2d}t\,K_2(\tan t)+\tan^{3d}t\,K_3(\tan t),\label{eq:schur}\\
 K_2&=FG+F^2/2,\qquad
 K_3=FG^2+F^2H+F^2G+F^3/3.\notag
\end{align}
This is an exact identity of Taylor coefficients.

We use the preceding all-order and eventual-$6.6m$ reports~\cite{AllOrders,Bound66}. They are bundled unchanged, with their reproducible sources. The needed facts are stated explicitly here.

First, $[t^N]P_m>0$ for every $m\ge2$ and even $d<N<3d$; moreover, for all sufficiently large $m$, this holds throughout $d<N\le3.3d$. Second, for even $3d\le N\le3.5d$,
\begin{equation}\label{eq:cubic}
 |[t^N]\tan^{3d}t\,K_3(\tan t)|
 <220(\log_2d+4)^2 2^d.
\end{equation}
The latter follows from the orbit marking $\binom{L+1}{2}$ for $H$, where $L$ is the deepest occupied dyadic factor. The complete proof is Lemma~2.1 of~\cite{Bound66}.

For the operator estimates, set
\begin{gather*}
 R=2/5,\quad w=49/100,\qquad
 \HH_a(x)=\prod_{j\ge1}\left(\cos\frac{\pi x}{2^j}+a\sin\frac{\pi x}{2^j}\right),\\
 h=\HH_R,\qquad L_R=h(1/2)=L(R).
\end{gather*}
On $X=\{f\in C^1(\mathbb R/\Z):f(0)=0\}$ use
\begin{gather*}
 N_0(f)=\sup_{x\notin\Z}\frac{|f(x)|}{t(x)h(t(x))^d},\qquad
 N_1(f)=\sup_x\frac{|f'(x)|}{h(t(x))^d},\\
 \|f\|_d=N_0(f)+\frac{(9/10)^d}{d}N_1(f),\qquad
 t(x)=\operatorname{dist}(x,\Z).
\end{gather*}
For even $d\ge128$ and $|a|\le R$, the inverse of $I-QV_a|_X$ has norm at most five. Evaluation at $1/2$ costs $L_R^d/2$. Also $1\le h\le7/5$ on $[0,1/2]$, $h$ increases on $[0,1/4]$, and each individual inverse-branch ratio satisfies
\begin{equation}\label{eq:branch}
 |W_a(y)|h(t(y))/h(t(x))\le1,
 \qquad W_a(y)=\cos\pi y+a\sin\pi y.
\end{equation}
The normalized frozen function $f=h_0+(I-QV_a)^{-1}q_a$ has the approximation
\begin{equation}\label{eq:repair}
 \begin{aligned}
 f&=\phi-C+e,\qquad \phi(x)=\HH_a(x)^d+\HH_a(x-1)^d,\\
 \|C\|_\infty,\|C'\|_\infty&\le610d(RL_R)^d,\qquad
 \|e\|_d\le60000d^2w^d.
 \end{aligned}
\end{equation}
Here $C$ is the cubic Hermite correction of the excess endpoint values and derivatives. Its scalar consequence is
\begin{equation}\label{eq:scalarerror}
 |F(a)-L(a)^d-L(-a)^d|\le31000d^2(wL_R)^d,
 \qquad wL_R<\beta:=11907/20000.
\end{equation}
The same estimates give the full-disk bounds
\begin{equation}\label{eq:fullFG}
 |F(a)|\le3L_R^d,\qquad |G(a)|\le4275dL_R^d.
\end{equation}
The proof of the $G$ bound uses $\|f-h_0\|_d\le1710d$ and the unchanged Green inverse; it does not extend an unregularized orbit sum beyond its convergence range.

Finally, writing $M_{d,q}=[a^q]B(a)^d$, the dyadic tangent comparison and low-degree orbit formula give
\begin{equation}\label{eq:lowbounds}
 |F_q|\le\tfrac52c^dM_{d,q},\qquad
 |G_q|\le(5\log_2d+18)c^dM_{d,q}\quad(0\le q\le d-2).
\end{equation}
A monomial carries weight one for $F$ and its deepest occupied index for $G$. Odd coefficients vanish. The tangent comparison and orbit identities are included in the preceding source archive.

\section{The actual quadratic asymptotic}
Write
\[
 A_k(d,N)=[t^N]\GG(t)^{kd}.
\]
\begin{proposition}\label{prop:quadratic}
Uniformly for even integers $N$ with $N/d$ in any compact subinterval of
$(2,2\mu(\arctan R))$,
\begin{equation}\label{eq:K2asymp}
 [t^N]\tan^{2d}t\,K_2(\tan t)
 =(2\log_2d+O(1))A_2(d,N).
\end{equation}
\end{proposition}
\begin{proof}
We first quantify the deepest occupied index. For $q/d$ in a compact subinterval of $(0,\infty)$, choose $a>0$ so that
\[
 \nu(a):=\sum_j\frac{ab_j}{1+ab_j}=q/d.
\]
Take $d$ independent labeled copies of each Bernoulli selector with probability $ab_j/(1+ab_j)$. Conditioning their total count on $q$ gives precisely the coefficient measure of $M_{d,q}$. Let $L$ be the largest selected index. Uniformly in this parameter range,
\begin{equation}\label{eq:depth}
 \mathbb E L=\log_2d+O(1).
\end{equation}
Here is a direct uniform justification. The total count has mass at its integer mean at least $c_0/\sqrt d$, by the lattice local limit theorem. All parameters remain in a compact positive interval; there are common exponential moments and a uniform span-one gap. The maximum mass after removing one selector, or after retaining only indices up to $J\ge1$, is at most $C_0/\sqrt d$: an independent binomial summand from the $j=1$ selectors supplies this bound, and convolution cannot increase it. The selector tail is comparable to $2^{-J}$. Marking one tail selector and, in the other direction, requiring all tail selectors to vanish, give
\[
 \mathbb P(L>J\mid q)\le\min(1,Cd2^{-J}),\qquad
 \mathbb P(L\le J\mid q)\le C e^{-cd2^{-J}}.
\]
Summation above and below $\lfloor\log_2d\rfloor$ proves~\eqref{eq:depth}, uniformly.

For even integers $q$ with $q/d$ in a compact subinterval of $(0,1)$, the two nearest half-integer orbits consequently give
\begin{equation}\label{eq:lowasymp}
 F_q=2c^dM_{d,q}(1+O(e^{-c_1d})),\qquad
 G_q=2c^dM_{d,q}(\log_2d+O(1)).
\end{equation}
Indeed the remote odd branches have an extra factor $n^{q-d}$, with at most logarithmic depth factors; their sums for $n\ge3$ are exponentially small uniformly when $q\le(1-\epsilon)d$.

Let $F_<,G_<$ retain raw degrees below $d$. From~\eqref{eq:fullFG}, their high tails on $|a|\le a_0<R$ are bounded by
\[
 |F-F_<|\le\frac{3L_R^d(a_0/R)^d}{1-a_0/R},\qquad
 |G-G_<|\le\frac{4275dL_R^d(a_0/R)^d}{1-a_0/R}.
\]
The low parts are bounded by a logarithmic factor times $L(a_0)^d$, using~\eqref{eq:lowbounds}. Their relative tail factor is
\[
 \tau(a_0)=\frac{a_0L_R}{RL(a_0)}<1.
\]
This follows because $a/L(a)$ increases on $(0,R]$:
$\nu(a)<a\sum b_j<2a\le4/5$. At the second saddle, $a_0=\tan t_2$ stays in a compact subinterval of $(0,R)$. Cauchy's estimate and the local-limit lower bound for $A_2$ therefore make every product containing a high tail exponentially negligible.

In the remaining low convolution, mark the two raw selector counts by $u,v$ in
\[
 c^{2d}\tan^{2d}t\,B(u\tan t)^d B(v\tan t)^d.
\]
At the positive saddle each count has mean $d\nu(a_0)$, lying in a compact subinterval of $(0,1)$. Conditional on total $t$-degree $N$, either count outside a fixed smaller interior interval has exponentially small probability. This follows by marking with $e^{\theta}$ and a Chernoff bound at the same saddle; division by the local-limit coefficient costs only $O(\sqrt d)$. The absolute bounds~\eqref{eq:lowbounds} control these atypical terms. Thus~\eqref{eq:lowasymp} applies on the whole convolution up to an $O(A_2)$ error.

Set $F_{\rm near}(a)=c^d(B(a)^d+B(-a)^d)$. Restoring its full nearest product instead of its truncation adds only the same atypical raw-degree tail. We obtain
\[
 [t^N]\tan^{2d}t\,F_<G_<
 =\log_2d\,[t^N]\tan^{2d}t\,F_{\rm near}(\tan t)^2+O(A_2).
\]
For even $d,N$, the last coefficient equals
\[
 2A_2(d,N)+2[t^N]\{\GG(t)\GG(-t)\}^d.
\]
The mixed term is exponentially negligible. On $|z|=t_2$ its product modulus is strictly below $\GG(t_2)^2$, uniformly on the compact saddle interval. Equality for the first factor would require $z/t_2=1$, while equality for the second would require $z/t_2=-1$, since consecutive coefficients are positive. Finally $F_<^2/2$ contributes only $O(A_2)$. This proves~\eqref{eq:K2asymp}.
\end{proof}

\section{A bounded local frozen correction}
The first-model saddle lies outside the Green disk. We do not approximate the actual frozen function there. Instead, subtract its explicit nearest model and estimate the remainder on the smaller second-saddle contour.

Let
\[
 I=[3.3,3.332],\quad
 t_-=.3413848,\quad t_+=.3484140,\quad
 \theta_0=153\pi/1000.
\]
The rational saddle certificate gives $t_2(\sigma)\in[t_-,t_+]$ for $\sigma\in I$. On the compact local arcs
\begin{equation}\label{eq:localset}
 a=\tan(te^{i\theta}),\qquad t_-\le t\le t_+,\quad |\theta|\le\theta_0,
\end{equation}
the exact local certificate proves, with strict slack,
\begin{equation}\label{eq:localconditions}
 \Re a>.303,\qquad .348<|a|<.38,\qquad
 |Q_a(2^{-j})|<1\quad(1\le j\le5),
 \quad Q_a(x)=\frac{\HH_a(1+x)}{\HH_a(1)\HH_a(x)}.
\end{equation}
A common strict upper bound below one is supplied for the five ratios. The certificate and its interval construction are described in Section~\ref{sec:certificate}.

Put $L_a=\HH_a(1/2)=L(a)$ and $g_a=aL_a$. The following lemma is the local analytic ingredient.
\begin{lemma}\label{lem:localfrozen}
Uniformly on the arcs~\eqref{eq:localset},
\begin{equation}\label{eq:localfrozen}
 F(a)=L(a)^d+L(-a)^d
       -a^dL(a)^{2d}\{\Psi_d(a)+O(d^{-1})\},
\end{equation}
where $\Psi_d$ is analytic on a common neighborhood and uniformly bounded there. In particular,
\[
 |F(a)-L(a)^d-L(-a)^d|\le C|aL(a)^2|^d.
\]
\end{lemma}
\begin{proof}
We give the principal-branch identity and all error estimates. The atomic function has the periodization
\[
 h_0(x)=\sum_{k\in\Z}\left(\frac{\sin\pi(x+k)}{\pi(x+k)}\right)^d.
\]
To distinguish its nonperiodic summand, write $s(x)=\sin(\pi x)/(\pi x)$. The frozen equation is
$f=V_af-a^dFh_0$. Divide by $\HH_a(x)^d$ and use
$\HH_a(x)=W_a(x/2)\HH_a(x/2)$. If $x_n=2^{-n-1}$, telescoping toward zero gives the exact, convergent identity
\begin{equation}\label{eq:principal}
 \frac{F}{L_a^d}=1+\sum_{n\ge0}
 \frac{W_a((1+x_n)/2)^d f((1+x_n)/2)-a^dFh_0(x_n)}{\HH_a(x_n)^d}.
\end{equation}
Each numerator vanishes at zero before it is summed.

The weighted error $f-\phi$ and its derivative are bounded by
$\operatorname{poly}(d)w^d h(t(x))^d$. For the derivative this needs a small refinement of~\eqref{eq:repair}, rather than division by its exponentially small norm weight. The two-row Green estimate gives
\[
 N_1(e)\le5dN_0(e)+\beta_dN_1(e)+N_1(\text{source}),
 \qquad \beta_d<3/5,
\]
with $N_0(e)=O(d^2w^d)$ and $N_1(\text{source})=O(d^2w^d)$. Hence $N_1(e)=O(d^3w^d)$, and the Hermite correction has the same exponential base.

For $0\le x\le1/2$, factorwise absolute values give
$|\HH_a(x)|\ge\HH_{\Re a}(x)\ge1$. The last inequality follows from real logarithmic concavity, the value one at zero, and the certified $L(.303)>1.05$. Thus $|g_a|>.348\cdot1.05=.3654$. The squared secondary weight
$|a\cos(\pi x/2)-\sin(\pi x/2)|^2$ has its maximum at an endpoint; its derivative changes sign only from negative to positive. By~\eqref{eq:localconditions}, its endpoint bounds, multiplied by $h\le7/5$, are below $.73$. Substituting $\phi$ for $f$ in the secondary branch of~\eqref{eq:principal}, and $F_{\rm near}=L(a)^d+L(-a)^d$ for $F$, therefore has normalized error
\[
 O(\operatorname{poly}(d)(.73w)^d),\qquad .73w=.3577<|g_a|.
\]
The rank-one error has smaller base $|a|\beta$. These error numerators still vanish at zero. Their normalized derivatives obey the same bound up to a polynomial factor, so summing the geometric sequence $x_n$ is legitimate.

The substituted secondary product is exactly
\[
 W_a((1+x)/2)^d\phi((1+x)/2)
 =\HH_a(1+x)^d+\HH_a(x-1)^d.
\]
Replacing $h_0(x)$ by $s(x)^d$ has exponentially smaller cost: the remaining periodization sum is $O((2/\pi)^d)$, with derivative $O(d(2/\pi)^d)$. Their vanishing at zero again controls the dyadic sum.

The negative-branch term at $n=0$ equals $L(-a)^d/L(a)^d$, recovering the second term of $F_{\rm near}$. Its counterterm is retained in the exponentially small remainder. For $n\ge1$, $0\le x_n\le1/4$ and
$|\HH_a(x_n-1)|\le|a|$: the first factor satisfies this because
$\tan(\pi/8)<2\Re a/(1-|a|^2)$, and every later factor has modulus at most one. Also $|L(-a)|<1$. The difference at zero vanishes, and its derivative has base $|a|^d$, exponentially smaller than $|g_a|^d$. The removed counterterm has the still smaller relative base $c|L(-a)|/|L(a)|^2<1$.

It remains to analyze the positive branch. Put
\[
 v_a(x)=\frac{s(x)}{\HH_a(x)},\qquad
 J(a)=\pi a-\frac{\HH_a'(1)}{\HH_a(1)}.
\]
The preceding estimates reduce~\eqref{eq:principal} to
\begin{equation}\label{eq:dyadicreduction}
 \frac{F_{\rm near}-F}{a^dL(a)^{2d}}
 =\sum_{n\ge0}\{v_a(x_n)^d-Q_a(x_n)^d\}+O(e^{-c_1d}).
\end{equation}
All constants are uniform on a neighborhood of the compact arcs, by the strict certificate margins.

We have $\Re J(a)>\pi\Re a>.9$. Indeed, for real $u\ge0$,
\[
 \Re\frac{a-u}{1+au}\le\Re a,
\]
because $\Re((1+a^2)/(1+au))>0$ when $\Re a>0$ and $|a|<1$. Summing after the first spatial factor gives
$\Re(\HH_a'(1)/\HH_a(1))\le\pi(\Re a-\Re(1/a))/2<0$.
For $0\le x\le1/64$, the real part of the first factor of $\HH_a(1+x)$ exceeds $.27$; the other real parts exceed $2/3$. All factor real parts of $\HH_a(x)$ exceed $.99$. Therefore
\[
 |(\log Q_a)''(x)|
 \le\pi^2(1+.38^2)
 \left\{\frac1{4(.27)^2}+\frac3{16}+\frac1{3(.99)^2}\right\}<46.
\]
Since $(\log Q_a)'(0)=-J(a)$, this proves $|Q_a(x)|\le e^{-x/2}$ on that interval. The five larger dyadic arguments are covered by~\eqref{eq:localconditions}. Also
\[
 v_a(x)=\prod_{j\ge1}(1+a\tan(\pi x/2^j))^{-1}
\]
has a uniform bound $|v_a(x)|\le e^{-c_2x}$ on $[0,1/2]$.

Uniform Taylor expansion now gives
\[
 \log v_a(x)=-\pi ax+O(x^2),\qquad
 \log Q_a(x)=-J(a)x+O(x^2).
\]
Replacing the powers in~\eqref{eq:dyadicreduction} by these exponentials costs $O(1/d)$: each error is at most $Cd x_n^2e^{-c_3dx_n}$, and the sum of $(dx_n)^2e^{-c_3dx_n}$ is uniformly bounded. The finitely many larger arguments cost exponentially little. Extending the resulting sum to all integer indices also costs exponentially little. Thus~\eqref{eq:localfrozen} holds with
\begin{equation}\label{eq:Psi}
 \Psi_d(a)=\sum_{n\in\Z}
 \left\{e^{-\pi ad2^{-n-1}}-e^{-J(a)d2^{-n-1}}\right\}.
\end{equation}
The positive real parts of $\pi a$ and $J(a)$ make this series uniformly convergent and bounded, with any fixed number of derivatives, on a common neighborhood. This completes the proof.
\end{proof}
For real $a$ in the corresponding interval, $J(a)>\pi a$ and $\Psi_d(a)>0$. Moreover $\Psi_{2d}(a)=\Psi_d(a)$, and its mean over one period in $\log_2d$ is $\log_2(J(a)/(\pi a))$, by Frullani's integral. Only its uniform boundedness is needed for the first-negative law.

\section{The outer contour and its exact certificate}\label{sec:certificate}
For $0<|a|\le R$, define
\[
 \omega(a)=\sup_{1\le X\le2}
 \frac{\max(|\HH_a(X)|,|\HH_{-a}(X)|)}{h(\operatorname{dist}(X,\Z))}.
\]
\begin{lemma}\label{lem:omega}
On any compact subset of $0<|a|\le R$,
\begin{equation}\label{eq:omega}
 |F(a)-L(a)^d-L(-a)^d|\le C d^2\{L_R\omega(a)\}^d.
\end{equation}
\end{lemma}
\begin{proof}
The continuous function $\omega$ has a positive lower bound on the compact set: its endpoint includes $|aL(a)|$, and $L(a)=cB(a)$ has no zero for $|a|<1$. The nearest-pair excess endpoint values are bounded by $\omega^d$, and its excess endpoint derivatives by $120d\omega^{d-1}$. Thus the fixed Hermite basis gives
$\|C\|_\infty+\|C'\|_\infty\le C_1d\omega^d$.

For $e=f-(\phi-C)$ the periodic vanishing source is
\[
 s=\HH_a(x-2)^d+\HH_a(x+1)^d-(V_aC-C)-a^d(\phi-C)(1/2)h_0,
 \qquad s(0)=0.
\]
The outer derivative satisfies
\[
 \frac{|(\HH_a(X)^d)'|}{h(\operatorname{dist}(X,\Z))^d}
 \le120d\omega^{d-1}\le C_2d\omega^d.
\]
For the transfer of $C$,~\eqref{eq:branch} and $h\ge1$ imply both
$|W_a(y)|^d/h(t(x))^d\le1$ and
$|W_a(y)|^{d-1}/h(t(x))^d\le1$.
Differentiation therefore gives $N_1(V_aC)\le C_3d^2\omega^d$.
For the rank-one term, $|a|^d|\HH_a(\pm1/2)|^d\le\omega^d$ by the endpoint identity; the correction involving $C$ is no larger. Since $\|h_0'\|\le\pi d$, its derivative has the same bound. Hence $N_1(s)\le C_4d^2\omega^d$.

For $t(x)\le1/4$, integrate from the nearest integer using $s(0)=0$ and increasing $h$. For $t(x)\ge1/4$, bound the weighted amplitudes directly and divide by $t(x)$. This yields $\|s\|_d\le C_5d^2\omega^d$. Apply the norm-five Green inverse, evaluate at $1/2$, and add the Hermite value correction. This proves~\eqref{eq:omega} with its exact exponential base.
\end{proof}

The outer certificate proves, with $\eta=1999/2000$, that
\begin{equation}\label{eq:outer}
 |a|L_R\frac{\max(|\HH_a(X)|,|\HH_{-a}(X)|)}{h(\operatorname{dist}(X,\Z))}
 <\eta\,g(a_0)^2,
\end{equation}
for
\[
 t_-\le t\le t_+,\quad \theta_0\le\theta\le\pi/2,\quad
 1\le X\le2,\qquad a=\tan(te^{i\theta}),\quad a_0=\tan t.
\]
Conjugation and evenness cover the rest of the circle.

For completeness, the certificate is a finite rational inequality proof, not a floating-point search. The complex tangent is enclosed by
\[
 \tan(u+iv)=\frac{\sin(2u)+i\sinh(2v)}{\cos(2u)+\cosh(2v)},
 \quad u=t\cos\theta,\ v=t\sin\theta.
\]
Machin's formula and Taylor remainder bounds provide outward integer intervals for all trigonometric and hyperbolic functions. Twenty spatial factors retain the same complex interval $a$. For $c_j=\cos(\pi X/2^j)$ and $s_j=\sin(\pi X/2^j)$, the exact squared factor is
\[
 |c_j\pm as_j|^2=c_j^2\pm2\Re(a)c_js_j+|a|^2s_j^2.
\]
Its logarithmic derivative is
\[
 2\frac{\pi}{2^j}
 \frac{\pm\Re(a)(c_j^2-s_j^2)+(|a|^2-1)c_js_j}{|c_j\pm as_j|^2}.
\]
For the omitted product, set $E=R\pi X_{\max}2^{-20}+(\pi X_{\max})^2/(6\cdot4^{20})$. The sum of the omitted factor deviations is at most $E<1$, so its modulus is at most $(1-E)^{-1}$. The omitted squared logarithmic derivative is bounded by
\[
 \frac{2\pi}{2^{20}}\frac{R+\tau}{1-R\tau-\tau^2/2},
 \qquad \tau=\frac{\pi X_{\max}}{2^{21}}.
\]
Spatial cells are split at $X=3/2$. For the real denominator, positivity and logarithmic concavity make its minimum on a distance cell occur at an endpoint. A midpoint value and an interval bound for the logarithmic derivative control a spatial box by $e^v\le1/(1-v)$ when $0\le v<1$. A direct squared-factor bound is used whenever that sharpening is unavailable.

The compact outer partition has 2812 boxes. Every squared ratio is below $999/1000<(1999/2000)^2$; its largest certified upper bound is approximately $.9989983569$. The local partition has 127 boxes and largest certified shallow $Q$ squared upper bound
\[
 \frac{99983441723839698176893318289531}{10^{32}}<1.
\]
The box partitions, exact rational endpoints, interval programs, and independent replay checks are included. The algorithm starts from the whole domain and bisects a box into two children; replay reconstructs the partition and checks every leaf. All new scalar constants, saddle coverage, and the fact that both model exponential bases exceed three on $I$ are also checked exactly. Strict margins give the compact-neighborhood slack used in Lemma~\ref{lem:localfrozen}.

\section{The coefficient comparison and the first negative degree}
\begin{proposition}\label{prop:comparison}
Uniformly for even $N$ with $N/d\in I$,
\begin{equation}\label{eq:comparison}
 [t^N]P_m(t)=2A_1(d,N)-(2\log_2d+O(1))A_2(d,N).
\end{equation}
\end{proposition}
\begin{proof}
Let $t=t_2(N/d)$. On the positive local arc, Lemma~\ref{lem:localfrozen} bounds the composed frozen remainder by $C|\GG(z)|^{2d}$. For the normalized positive coefficient law at radius $t$, with probabilities $p_j$, one has
\[
 \left|\sum_jp_je^{ij\theta}\right|^2
 \le1-2p_1p_2(1-\cos\theta).
\]
The two probabilities are uniformly positive, so this gives a uniform Gaussian modulus bound. Integration over the arc retains the factor $d^{-1/2}$. The negative local arc has the same bound by evenness. On the outer arcs, Lemma~\ref{lem:omega} and~\eqref{eq:outer} give an exponentially smaller term, even after division by the second-model local-limit coefficient. Thus
\[
 [t^N]\tan^dt\,F(\tan t)=2A_1(d,N)+O(A_2(d,N)).
\]
The exact factor two comes from the equal even coefficients of $\GG(t)^d$ and $\GG(-t)^d$.

Proposition~\ref{prop:quadratic} handles the quadratic term. By~\eqref{eq:cubic}, the cubic contribution is polynomial in $\log d$ times $2^d$. Both model exponential bases exceed three uniformly on $I$, as certified, so this is exponentially negligible relative to $A_2$. The cosine term is $O(d)$ by Cauchy's estimate on $|t|=1$. The exact expansion~\eqref{eq:schur}, valid since $N<4d$, proves~\eqref{eq:comparison}.
\end{proof}

\begin{proof}[Proof of Theorem~\ref{thm:main}]
Uniform positive-series saddle asymptotics give
\begin{equation}\label{eq:ratio}
 \frac{A_2(d,N)}{A_1(d,N)}
 =C(\sigma)e^{d\Delta(\sigma)}(1+o(1)),\qquad
 C(\sigma)=\sqrt{\frac{v_1(\sigma)}{2v_2(\sigma)}},
 \quad v_k=t_k\mu'(t_k),\quad \sigma=N/d.
\end{equation}
The errors are uniform on $I$. The local-limit proof uses a common analytic neighborhood, positive variance bounds, and the uniform span-one gap from consecutive positive coefficients. In particular $C$ is bounded above and below by positive constants.

Since $\Delta$ increases and vanishes at $\sigma_*$,~\eqref{eq:comparison} is positive uniformly below $\sigma_*-\epsilon$ and negative at $\sigma_*+\epsilon$, for each sufficiently small fixed $\epsilon>0$ and all large $d$. The preceding eventual-$6.6m$ theorem covers every earlier degree to the left of $I$. Therefore $N_m/d\to\sigma_*$.

For the sharper location, put
\[
 M_d=d\sigma_*-\frac{\log\log d}{\Lambda}.
\]
Uniformly when $N-M_d$ is bounded,
\[
 d\Delta(N/d)=-\log\log d+\Lambda(N-M_d)+o(1).
\]
Choose constants $0<C_{\min}\le C(\sigma)\le C_{\max}$ on $I$. Since $\Delta$ increases, every degree below $M_d-K$ has
\[
 \log_2d\,\frac{A_2}{A_1}
 \le\frac{C_{\max}}{\log2}e^{-\Lambda K+o(1)}.
\]
A sufficiently large fixed $K$ makes this less than $1/2$, proving positivity at all such degrees within $I$. The earlier theorem covers the rest. At an even degree in $[M_d+K,M_d+K+2]$, the corresponding lower bound using $C_{\min}$ exceeds two, giving a negative coefficient. No monotonicity of the full coefficient ratio is assumed. We conclude $N_m=M_d+O(1)$, which is~\eqref{eq:main}.
\end{proof}

\section{A refined lattice window}
Let $v_k=t_k\mu'(t_k)$ at the crossing, and put
\[
 C_*=\sqrt{v_1/(2v_2)},\qquad
 u_d=d\sigma_*-\frac{\log\log d}{\Lambda}
                +\frac{\log((\log2)/C_*)}{\Lambda}.
\]
\begin{proposition}\label{prop:rounding}
There are constants $C,B,d_0$ such that, for even $d\ge d_0$, every even post-cancellation degree $N<u_d-C/\log d$ has positive pressure coefficient, while every even degree with
\[
 u_d+C/\log d<N\le u_d+B
\]
has negative coefficient. In particular, whenever
$\operatorname{dist}(u_d,2\Z)>C/\log d$,
\[
 N_m=2\left\lceil\frac{u_d}{2}\right\rceil.
\]
\end{proposition}
\begin{proof}
The compact family of positive saddle laws gives the quantitative version of~\eqref{eq:ratio}, with relative error $O(d^{-1/2})$. Indeed, on a fixed sufficiently small central arc, the third-order characteristic expansion bounds the integral error by
\[
 O\left(d\int |\theta|^3e^{-cd\theta^2}\,d\theta\right)=O(d^{-1}),
\]
against a main integral of order $d^{-1/2}$. The complement has a uniform exponential gap. This makes the error uniform in the compact slope interval.

For fixed $B$ and $|N-u_d|\le B$, Taylor expansion at $\sigma_*$ and~\eqref{eq:comparison} therefore give
\[
 \frac{[t^N]P_m(t)}{2A_1(d,N)}
 =1-e^{\Lambda(N-u_d)}+O_B(1/\log d).
\]
The all-earlier-sign argument already gives positivity below $u_d-B_0$ for some fixed $B_0$. Choose $B>\max(B_0,3)$. On $[-B,B]$, the derivative of $e^{\Lambda x}$ is bounded below by a positive constant. A sufficiently large fixed $C$ gives the two stated sign bounds. If the center is farther than $C/\log d$ from every even integer, the next even integer is negative and every earlier post-cancellation degree is positive, proving the last assertion.
\end{proof}
The constants are existential. No claim is made about how often $u_d$ approaches the even lattice, so this is not an unconditional exact floor rule.

\section{Scope and reproducibility}
The first-negative asymptotic is proved by the analytic comparison and exact compact-domain certificates. It is not inferred from a finite fit. Its slope is strictly below $20/3$, so a floor rule with that limiting slope cannot hold eventually. No unconditional rounding formula or effective onset threshold is supplied.

The real frozen correction contains a bounded oscillation in $\log_2d$. Its contribution is smaller than the quadratic logarithmic multiplier, which is why it does not alter the leading slope or logarithmic correction. Proposition~\ref{prop:rounding} isolates an $O(1/\log d)$ ambiguity window around the even lattice; it does not determine how often that window is encountered.

The source archive contains the compact local and outer partitions, portable exact checkers, model-crossing certificate, source hashes, and unchanged analytic prerequisites. The supplied first-negative finite data, when distributed as a companion, are supplementary and are not a premise of the theorem.

\begin{thebibliography}{9}
\bibitem{Repo} V.~Reshetnikov, \emph{Thue--Morse Integer Pressure}, repository research manuscript, pinned at commit \texttt{63a7a325109ba611a1816b61dfd0eb072b896a7a}. Source included.\newline
\url{https://github.com/VladimirReshetnikov/ProveIt/blob/63a7a325109ba611a1816b61dfd0eb072b896a7a/Analysis/FabiusFunction/docs/semi-formalized-research-frontiers/drafts/thue-morse/Thue_Morse_Integer_Pressure/article.tex}
\bibitem{AllOrders} \emph{The Full Pressure Positivity Range at Every Integer Order}, research report prepared for V.~Reshetnikov with OpenAI, 1 October 2026. Unchanged PDF and reproducible sources included.
\bibitem{Bound66} \emph{An Asymptotic Lower Bound for the First Negative Thue--Morse Pressure Degree}, research report prepared for V.~Reshetnikov with OpenAI, 1 October 2026. Unchanged PDF and TeX included; supplies the earlier-degree window and cubic envelope.
\bibitem{Signs} \emph{Infinite Taylor Sign Changes in Integer Thue--Morse Pressure}, research report prepared for V.~Reshetnikov with OpenAI, 1 October 2026. Unchanged PDF included.
\bibitem{FSS} A.~Fan, J.~Schmeling and W.~Shen, \emph{Multifractal analysis of generalized Thue--Morse trigonometric polynomials}, arXiv:2212.13234v1 (2022). \url{https://arxiv.org/html/2212.13234v1}
\bibitem{BGKS} M.~Baake, P.~Gohlke, M.~Kesseb\"ohmer and T.~Schindler, \emph{Scaling properties of the Thue--Morse measure}, arXiv:1810.06949v2 (2019). \url{https://arxiv.org/html/1810.06949v2}
\end{thebibliography}
\end{document}
